% !PRINTED 463
\weberpage{463}
%<<CONT
\[
\Phi(t) = (t - \Theta)(t - \Theta_1) \ldots (t - \Theta_{n-1}),
\]
then
\[
\Phi(t) \left( \frac{\omega}{t - \Theta} + \frac{\omega_1}{t - \Theta_1} + \cdots + \frac{\omega_{n-1}}{t - \Theta_{n-1}} \right) = \Psi(t)
\]
is an integral rational function of degree $n - 1$ of $t$,
whose coefficients are numbers in $\Omega$, and from this it
follows, on setting $t = \Theta$,
\[
\omega = \frac{\Psi(\Theta)}{\Phi'(\Theta)},
\]
in which $\Phi'(\Theta)$ is different from zero.

\gesperrt{Accordingly the field $\Omega(\Theta)$ is identical with the field
$\Omega(\alpha)$.}

\gesperrt{4. If $\Theta$ is not primitive, then not every quantity
in $\Omega(\alpha)$ can be expressed rationally in terms of $\Theta$. The
field $\Omega(\Theta)$ is a divisor of the field $\Omega(\alpha)$,
and the degree of $\Omega(\Theta)$ is a divisor of the degree of
$\Omega(\alpha)$.}

For every number of the field $\Omega(\Theta)$ satisfies an equation
in $\Omega$ of degree $n_1$, where $n_1$ is a divisor of $n$ and
less than $n$. Hence a primitive quantity of the field
$\Omega(\alpha)$, which satisfies an irreducible equation of degree
$n$, cannot be contained in $\Omega(\Theta)$.

\gesperrt{The field $\Omega(\alpha)$ is called primitive if, apart from the
quantities in $\Omega$, it contains no imprimitive quantities; imprimitive,
if it contains still other imprimitive quantities.}

From this definition it follows first of all that a field whose
degree is a prime is necessarily primitive; for an imprimitive
quantity $\Theta$ in $\Omega(\alpha)$ satisfies an equation whose
degree is a divisor of the degree $n$ of $\Omega(\alpha)$ different
from $n$. If this divisor is equal to 1, then $\Theta$ is contained
in $\Omega$.

We shall now become acquainted with some of the most important
properties of the imprimitive fields.

If $\Omega(\alpha)$ has a divisor $\Omega'$ different from
$\Omega$, which in its turn contains $\Omega$ as a divisor (so that
$\Omega'$ is a field over $\Omega$), and if $\beta$ is a quantity that
belongs to the field $\Omega'$ but not to $\Omega$, then the
field $\Omega(\beta)$ is an algebraic field over $\Omega$ and
at the same time a divisor of $\Omega'$ and of $\Omega(\alpha)$, and
according to what we have just
%CONT>>
